\section{Security and Privacy Analysis}
\label{sec:security}

This section surveys the security and privacy risks associated with Computer-Use Agents (CUAs). Relative to conventional LLM safety, which has largely focused on dialogue behavior and text generation, CUA security must additionally account for \emph{action-level} risk: these systems perceive visual environments, interpret user goals, and execute low-level operations over external interfaces, thereby exposing vulnerabilities at the intersection of perception, reasoning, and execution~\cite{jones2025systematization,tian2023evil}. 

Following the lifecycle perspective introduced in Section~\ref{sec:lifecycle}, we organize the discussion according to the stages at which security and privacy failures are introduced, amplified, or realized. We first specify a threat model, then introduce an attribution framework that distinguishes model-subjective factors (curiosity and malice) from non-subjective environmental causes, and finally review representative threats across creation, deployment, execution, and maintenance. Table~\ref{tab:cua_threats} summarizes this taxonomy by lifecycle stage and attribution class. Because the survey is motivated in part by open deployment settings exemplified by OpenClaw, we place particular emphasis on risks that emerge when CUA capabilities are exposed through persistent communication channels, extensible tools, and long-lived execution contexts.

\begin{table*}[!t]
\centering
\small
\caption{Threat taxonomy by lifecycle stage and attribution (C: curious, M: malicious, E: environmental). Abbreviations: PoT, Plan-of-Thought; MCP, Model Context Protocol; TOCTOU, time-of-check to time-of-use; IPI, Indirect Prompt Injection; EIA, Environmental Injection Attack; RAG, Retrieval-Augmented Generation; PII, personally identifiable information.}
\label{tab:cua_threats}
\setlength{\tabcolsep}{5pt}
\renewcommand{\arraystretch}{1.15}
\begin{tabular}{>{\raggedright\arraybackslash}p{1.55cm}
                >{\raggedright\arraybackslash}p{4.2cm}
                >{\raggedright\arraybackslash}p{1.55cm}
                >{\raggedright\arraybackslash}p{8.1cm}}
\toprule
\textbf{Stage} & \textbf{Threat Type} & \textbf{Attribution} & \textbf{Representative Work} \\
\midrule
Creation & PoT backdoor, training-time poisoning & M & BadVLA~\cite{zhou2025badvla}, Hidden Ghost Hand~\cite{cheng2025hidden} \\
Creation & Unsafe action interfaces and safety gaps & E & OS-Harm~\cite{kuntz2025harm}, CUAHarm~\cite{tian2025measuring} \\
Deployment & Over-privilege, weak isolation & E & Jones~\cite{jones2025systematization}, CSAgent~\cite{gong2025secure} \\
Deployment & Supply-chain risk (skills, MCP) & M or E & Les Dissonances~\cite{li2025dissonances}, MCP Security Bench~\cite{zhang2025mcp} \\
Execution & TOCTOU, action rebinding & E & Zero-Permission Manipulation (``Action Rebinding'')~\cite{qian2026zero}, Hijacking JARVIS (``AgentHazard'')~\cite{liu2025hijacking} \\
Execution & Clickjacking, visual anchors & E & Chameleon~\cite{zhang2026realistic}, WAInjectBench~\cite{liu2025wainjectbench} \\
Execution & IPI, EIA & M or C & InjecAgent~\cite{zhan2024injecagent}, WASP~\cite{evtimov2025wasp}, Chameleon~\cite{zhang2026realistic} \\
Execution & Memory/RAG poisoning & M & AgentPoison~\cite{chen2024agentpoison} \\
Execution & Curiosity-driven overreach, PII leakage & C or M & AgentDAM~\cite{zharmagambetov2025agentdam}, SafeArena~\cite{tur2025safearena} \\
Execution & Agent-to-agent propagation & C or M & Chen~\cite{chen2026openclaw}, Prompt Infection~\cite{lee2024prompt} \\
Maintenance & Governance gaps & E & Aegis~\cite{adapala2025aegis}, Agentic AI Survey~\cite{lazer2026survey} \\
\bottomrule
\end{tabular}
\end{table*}

\subsection{Threat Model}
\label{subsec:threat_model}

To structure the discussion, we adopt a threat-modeling perspective that identifies attacker roles, protected assets, trust boundaries, platform assumptions, and the major classes of threats emphasized in the recent literature.

\textbf{Attackers.} Prior work suggests at least three recurring attacker roles. First, \emph{external attackers} may control web pages, documents, search results, screenshots, or Model Context Protocol (MCP) tool returns later consumed by the agent. Second, \emph{supply-chain attackers} may distribute malicious or vulnerable skills, plugins, wrappers, or tool integrations. Third, \emph{malicious capability providers} may intentionally publish trojanized services or skills to agent ecosystems. These categories are not mutually exclusive, and some attacks combine multiple layers of compromise.

\textbf{Assets.} The primary assets include: (i) the \emph{confidentiality} of user data such as emails, credentials, files, and personally identifiable information; (ii) the \emph{integrity} of task execution, i.e., whether the agent’s actions remain faithful to user intent; and (iii) control over \emph{privilege-bearing interfaces}, including operating-system capabilities, online accounts, and external services. Compromise along any of these dimensions may result in privacy loss, unauthorized action, or persistent exposure.

\textbf{Trust boundaries.} A core challenge for CUAs is that user intent, environmental content, memory state, and tool-mediated outputs are often processed within a shared reasoning loop. The first trust boundary separates \emph{user-authored intent} from \emph{runtime content}, including interface text, web content, screenshots, retrieved snippets, and documents that may contain hidden or misleading instructions. The second separates \emph{trusted interfaces and tools} from \emph{untrusted external sources}, such as third-party web pages, documents, or MCP servers whose outputs may be incorrect, adversarial, or poisoned. When these boundaries are weakly represented or insufficiently enforced, spoofing, tampering, and privilege escalation become easier to realize.

\textbf{Platform scope.} Risk also varies across \emph{single-agent} and \emph{multi-agent} platforms. In multi-agent settings, prompts, capabilities, and learned behaviors may propagate across agents with limited human review, increasing the difficulty of attribution, containment, and audit. This consideration is especially relevant in open deployment settings, including OpenClaw-like ecosystems, where persistent communication channels, extensible tools, and agent-to-agent interaction can blur the distinction between local execution, external content, and socially propagated capability.

\textbf{Assumptions.} We assume that the underlying model can be induced, misdirected, or manipulated by external content; that the execution environment (web, desktop, or mobile) is only partially trusted; and that the available action space includes high-impact operations such as payment, deletion, installation, authentication, or command execution. Under these assumptions, relatively small failures in intent grounding or interface isolation may produce disproportionate downstream consequences.

\textbf{Threat types.} Following prior systems work, CUA-specific risks can be related to STRIDE-style categories. \emph{Spoofing} includes Indirect Prompt Injection (IPI) and Environmental Injection Attack (EIA), where attacker-controlled content is interpreted as instruction. \emph{Tampering} includes memory poisoning, Retrieval-Augmented Generation (RAG) poisoning, and training-time backdoors. \emph{Repudiation} concerns incomplete provenance, weak audit trails, and insufficient action justifications. \emph{Information disclosure} includes leakage of credentials, documents, and Personally Identifiable Information (PII) across tools or applications. \emph{Denial of service} includes hijacking and resource-exhaustion behaviors that prevent correct task completion. \emph{Elevation of privilege} includes over-privileged interfaces, weak isolation, and action rebinding.

From a lifecycle perspective, these threats manifest differently across stages. During \emph{Creation/Development}, adversaries primarily target training data, model internals, and action abstractions, yielding backdoors or unsafe interfaces. During \emph{Deployment}, missing provenance, weak isolation, or excessive permissions enlarge the practical attack surface. During \emph{Execution}, open and adversarial environments create opportunities for interface deception, prompt injection, and memory poisoning that directly affect confidentiality, integrity, and availability. During \emph{Maintenance}, continuous updates, capability sharing, and ecosystem expansion introduce propagation, regression, and governance challenges. The attribution framework below refines this view by distinguishing model-subjective from predominantly environmental causes of unsafe behavior.

\subsection{Attribution Framework: LLM Subjectivity}
\label{subsec:attribution}

A recurring issue in the agent-security literature is that unsafe behavior cannot always be attributed in the same way. To clarify this distinction, we adopt a simple attribution framework based on whether the model exhibits subjective harmful intent.

\textbf{Curiosity} refers to cases in which the model lacks harmful intent but nevertheless departs from the user’s intended scope. Examples include exploring irrelevant interfaces, collecting unnecessary information, or following salient but non-essential content. Such behavior may increase exposure to adversarial inputs and violate data-minimization expectations, but it does not imply that the model has internalized an attacker’s objective.

\textbf{Malice} refers to cases in which the model has been induced---for example through injection, poisoning, or backdoor triggers---to internalize and operationalize an adversarial goal. In this regime, the model does not merely fail passively; rather, it behaves as though it has adopted a harmful objective such as exfiltration, sabotage, or unauthorized redirection of task flow.

\textbf{Environmental} attribution refers to threats that arise without subjective fault on the part of the model. Representative cases include time-of-check to time-of-use (TOCTOU) mismatches, clickjacking, and other forms of perception--action misbinding caused by adversarial or unstable interface conditions. In such scenarios, the primary failure lies in the environment or system design rather than in the model’s intent.

This tripartite framing is intended as an analytic convenience rather than a claim about philosophical agency. Its purpose is to provide a practical vocabulary for surveying the literature. In the following subsections, each threat class is discussed in lifecycle context and annotated as curious, malicious, or environmental where appropriate.

\subsection{Threats in the Creation/Development Stage}
\label{subsec:threats_creation}

Threats introduced during creation or development originate in model design, training data, alignment procedures, or action-interface construction. Although such vulnerabilities may remain latent until deployment or execution, their root causes lie upstream in how the CUA is built. Existing studies mainly emphasize two themes: adversarial manipulation of model internals and unsafe exposure of action primitives.

\subsubsection{Plan-of-Thought Backdoors and Training-time Poisoning}

Within the creation stage, training-time poisoning and Plan-of-Thought (PoT) backdoors are most naturally attributed to \textbf{malice}. In these attacks, adversaries implant triggers into vision-language-action models such that apparently benign inputs induce harmful actions at inference time while nominal performance on clean inputs remains largely intact. BadVLA~\cite{zhou2025badvla}, for example, shows that poisoned training can redirect agent behavior toward attacker-specified endpoints with near-perfect success rates. Hidden Ghost Hand~\cite{cheng2025hidden} further demonstrates composite triggers for mobile GUI agents, reporting attack accuracies as high as 99.7\%. Agent Security Bench~\cite{zhang2024agent} formalizes PoT backdoors as a distinct agent-specific threat class. Taken together, this line of work suggests that harmful objectives can be embedded during model construction yet remain dormant until triggered later in execution.

\subsubsection{Unsafe Action Interfaces and Design Gaps}

A second creation-stage theme concerns \textbf{environmental} risk arising from unsafe action abstractions. Prior systems analyses argue that low-level, high-entropy action spaces are inherently harder to secure, whereas better interface abstractions can reduce the range of dangerous behaviors exposed to model output~\cite{jones2025systematization}. When high-impact operations such as payment confirmation, deletion, installation, or command execution are delegated too directly, the resulting system becomes more susceptible to spoofing, misexecution, and privilege escalation. Benchmarks such as OS-Harm~\cite{kuntz2025harm} and CUAHarm~\cite{tian2025measuring} further indicate that safety measurements are highly sensitive to interface design choices. From a survey perspective, these studies point to interface construction itself as part of the security problem, rather than merely a downstream implementation detail.

\subsection{Threats in the Deployment Stage}
\label{subsec:threats_deployment}

Deployment-stage threats emerge when CUAs are connected to real interfaces, tools, services, and user environments. In many cases, these threats do not create entirely new failure modes by themselves; instead, they amplify or enable downstream execution-time attacks by weakening provenance, isolation, or permission boundaries.

\subsubsection{Over-privilege and Weak Isolation}

The literature generally treats over-privilege and weak isolation as \textbf{environmental} deployment failures. Jones et al.~\cite{jones2025systematization} identify missing provenance, weak interaction binding, and unclear privilege boundaries as recurring system-level risks for computer-use agents. If provenance is missing, the agent cannot reliably distinguish user intent from attacker-controlled instruction. If the binding between perception and execution is weak, the system becomes vulnerable to clickjacking or action rebinding. If privilege boundaries are ambiguous, the agent may overreach into sensitive applications or accounts. These concerns appear especially acute in OpenClaw-like deployments, where persistent sessions, multi-channel ingress, and extensible tool bindings can concentrate heterogeneous authority behind a single user-facing endpoint. Across such settings, the literature consistently points to least privilege, sandboxing, and explicit trust separation as baseline deployment mechanisms rather than optional hardening features.

\subsubsection{Supply-chain Risk in Skills and MCP Tools}

Supply-chain risk in skills and MCP-integrated tools spans both \textbf{malicious} and \textbf{environmental} attribution. Malicious supply-chain risk arises when a provider intentionally distributes trojanized skills, hostile MCP servers, or tool definitions designed to exfiltrate data, manipulate downstream reasoning, or escalate privilege. MCP Security Bench~\cite{zhang2025mcp} systematizes twelve attack types across tool discovery, argument passing, and result return, while Les Dissonances~\cite{li2025dissonances} documents real-world tools that are capable of confidential-data harvesting or pollution by design.

Environmental supply-chain risk, by contrast, arises when tools are not intentionally malicious but nevertheless remain exploitable because of excessive permissions, inadequate isolation, or unsafe assumptions about tool outputs. Les Dissonances~\cite{li2025dissonances} evaluates 73 real-world LangChain and LlamaIndex tools and reports broad susceptibility to hijacking, cross-tool harvesting, and cross-tool pollution. One implication of this literature is that deployment-time integrity cannot be reduced to malware screening alone: secure composition also depends on capability scoping, provenance tracking, and the defensive treatment of tool outputs as potentially adversarial inputs.

\subsubsection{Insufficient Pre-release Red Teaming}

Another deployment-stage concern is the absence of realistic adversarial evaluation before release. This is most naturally viewed as an \textbf{environmental} gap in the deployment pipeline. RedTeamCUA~\cite{liao2025redteamcua}, for instance, argues that hybrid web--OS environments expose interaction paths that are not well captured by simplified benchmarks. Without such testing, systems may enter real deployment with vulnerabilities that are only discovered after exposure. From a lifecycle viewpoint, pre-release red teaming links creation and maintenance: it can reveal design flaws prior to release while also seeding regression suites and runtime monitoring rules for later stages.

\subsection{Threats in the Execution Stage}
\label{subsec:threats_execution}

The execution stage is where most CUA threats become operational. At this stage, models interact with open environments, user data, dynamic interfaces, and external tools under real-time constraints. Accordingly, the execution literature spans the full attribution spectrum, including environmental manipulation, curiosity-driven overreach, and malicious objective internalization.

\subsubsection{TOCTOU and Action Rebinding}

TOCTOU and action-rebinding attacks are typically classified as \textbf{environmental}. The central issue is that the agent’s perceptual state and its subsequent action may become desynchronized. Qian et al.~\cite{qian2026zero} report a ``reasoning gap'' of 4--15 seconds between screen understanding and action execution, during which an attacker can alter interface state and redirect the intended click toward a malicious target. Their study demonstrates 100\% authorization hijacking across six Android agents. Similarly, AgentHazard~\cite{liu2025hijacking} shows that deceptive overlays and misleading content can steer agents at non-trivial success rates. In this class of attack, the model’s initial intent may be correct, while the unsafe outcome is produced by environmental change between observation and actuation.

\subsubsection{Clickjacking and Visual Anchor Injection}

A related line of work studies \textbf{environmental} attacks on visual perception. Rather than manipulating reasoning directly, these attacks perturb the perceptual cues on which action selection depends. Chameleon~\cite{zhang2026realistic} shows that optimized visual triggers occupying only a small portion of the screen can substantially increase attack success rates. WAInjectBench~\cite{liu2025wainjectbench} further indicates that textual filtering is insufficient against imperceptible or cross-modal perturbations. Collectively, these studies suggest that CUA security cannot be understood solely as a prompt-level or text-level alignment problem; visual anchoring, layout semantics, and perception robustness are also part of the attack surface.

\subsubsection{Indirect Prompt Injection and Environmental Injection}

Indirect Prompt Injection (IPI) and Environmental Injection Attacks (EIA) occupy a central place in the execution-stage literature. Their attribution may be \textbf{malicious} when the model internalizes attacker-supplied objectives, or \textbf{curious} when it follows salient but irrelevant content without adopting an explicitly harmful goal. The common mechanism is that instructions embedded in web pages, retrieved documents, search results, screenshots, or MCP tool returns are processed as though they were authoritative. InjecAgent~\cite{zhan2024injecagent} benchmarks 1,054 cases and reports substantial vulnerability for ReAct-style prompting, with attack rates increasing further under hacking-oriented prompts. WASP~\cite{evtimov2025wasp} reports high deception rates in end-to-end web navigation, while SafeSearch~\cite{dong2025safesearch} shows that poisoned search results alone can yield a 90.5\% attack success rate. Chameleon~\cite{zhang2026realistic} and GhostEI-Bench~\cite{chen2025ghostei} extend this picture to dynamic and mobile environments. Across these studies, a recurring conclusion is that environmental content and tool outputs should not be treated as neutral context, but rather as potentially adversarial inputs~\cite{zhang2025browsesafe}.

\subsubsection{Memory and RAG Poisoning}

Memory poisoning and RAG poisoning are generally attributed to \textbf{malice}. Here the attacker seeks not merely to redirect a single episode, but to implant persistent harmful influence into the retrieval or memory substrate on which future behavior depends. AgentPoison~\cite{chen2024agentpoison} optimizes poisoned documents so that retrieval later activates latent malicious instructions during task execution. Agent Security Bench~\cite{zhang2024agent} similarly identifies memory poisoning as more persistent, and therefore potentially more severe, than one-shot injection. Les Dissonances~\cite{li2025dissonances} complements this perspective through its analysis of cross-tool harvesting and pollution. A broader synthesis from this literature is that the unit of compromise shifts from the immediate prompt to the longer-lived informational state that shapes future reasoning.

\subsubsection{Curiosity-driven Overreach and Data Minimization Violations}

Not all privacy failures arise from explicit adversarial compromise. A growing body of work instead highlights \textbf{curiosity-driven} overreach, in which agents exceed task scope, inspect irrelevant interfaces, or collect unnecessary data. AgentDAM~\cite{zharmagambetov2025agentdam} shows that mainstream agents often fail to satisfy data-minimization expectations and may capture or leak irrelevant PII. By contrast, a \textbf{malicious} interpretation is more appropriate when the model has been induced to exfiltrate or misuse data under attacker influence. SafeArena~\cite{tur2025safearena} finds that GPT-4o can complete a substantial fraction of harmful tasks on real websites, and prior systems work suggests that tool chaining may facilitate data movement across applications~\cite{jones2025systematization,li2025dissonances}. The literature therefore points to a dual requirement: task-boundary discipline is needed to limit benign overcollection, while adversarial robustness is needed to prevent attacker-directed exfiltration.

\subsubsection{Long-horizon Constraint Forgetting}

Long-horizon execution introduces additional failure modes that are partly \textbf{curious} and partly \textbf{environmental}. LPS-BENCH~\cite{chen2026lps} suggests that agents often lack sufficient foresight to reason about cascading consequences across multi-step plans, allowing goal completion pressure to dominate safety constraints. Related results from OS-Harm~\cite{kuntz2025harm} and CUAHarm~\cite{tian2025measuring} indicate that dialogue-level alignment does not automatically transfer to action-oriented settings. In aggregate, this literature points to a temporal dimension of CUA risk: safety constraints must remain stable not only within single decisions, but across extended interaction horizons in which partial progress can gradually erode earlier safeguards.

\subsubsection{Agent-to-agent Propagation and Cross-agent Threats}

Cross-agent interaction introduces a distinct execution-stage threat surface. Attribution again spans both \textbf{curiosity} and \textbf{malice}. On the non-malicious side, agents may exchange skills, demonstrations, or knowledge in ways that are locally useful but globally difficult to govern. Chen et al.~\cite{chen2026openclaw} provide the first large-scale empirical study of peer learning among AI agents on Moltbook, reporting participation by 2.4M agents, an 11.4:1 ratio of teaching to help-seeking statements, and machine-scale propagation of skills with minimal human review. These findings are especially relevant to open ecosystems such as OpenClaw, where persistent communication channels and socially mediated capability sharing can transform local mistakes or poisoned content into ecosystem-level exposure.

On the malicious side, Prompt Infection~\cite{lee2024prompt} shows that harmful prompts can replicate across agents in multi-agent systems, while Aegis~\cite{adapala2025aegis} and the Agentic AI Cybersecurity Survey~\cite{lazer2026survey} identify collusion, cascading failure, and oversight evasion as emerging platform-level concerns. Multi-Agent Security Tax~\cite{peigne2025multi} further highlights the trade-off between stronger security controls and collaborative flexibility. Taken together, these studies suggest that agent societies require security mechanisms beyond the single-agent paradigm, including identity verification, provenance-preserving exchange, semantic malware containment, and governance policies for capability diffusion.

\subsection{Threats in the Maintenance Stage}
\label{subsec:threats_maintenance}

Maintenance-stage threats concern long-term operation after release. They are primarily \textbf{environmental} in attribution, but remain central to whether CUA deployments stay safe under model updates, interface drift, evolving tools, and expanding ecosystems.

\subsubsection{Governance Gaps and Continuous Evaluation}

A consistent theme in recent work is that CUA security cannot be treated as a one-time property established before deployment. Model updates, tool changes, interface redesign, and policy adjustments may all introduce regressions. Benchmarks such as OS-Harm can serve as regression suites, while broader evaluation frameworks track shifts in harmful capability over time. Enterprise-oriented studies further emphasize auditability, incident response, provenance-aware logging, and governance controls~\cite{jones2025systematization,zhang2024agent}. In this sense, the maintenance problem is not merely one of operational hygiene; it is a structural requirement for CUA ecosystems in which capabilities, interfaces, and trust relationships continue to evolve after release. This observation directly motivates the lifecycle-aligned defenses and governance recommendations discussed in Section~\ref{sec:discussion}.